\section{The entire bounded-$h/q$ part of the doubly small endpoint}
\label{sec:new-c144-bounded-h-ratio}

\paragraph{Source scope and result.}
Retain the original entered physical C14.4 source, all strict gaps and
all original prior-complement intersections. Assume
\[
 S=\sigma+L,\ \kappa=x/S,\ \nu=\sigma/S
       \in\operatorname{Inf}_{>0},\qquad
 \ell=1-\nu,\quad q=\nu/\ell,\quad
 \tau=h/q\ \hbox{is bounded}.
 \tag{BN1}
\]
The complete original parent is a finite separated tensor SOS for
every test and every real or complex state on BN1. The quantity
$\beta/q^2$ is unrestricted, including zero and unbounded values.
The proof below is a fresh circuit completion of SN, not an import of
a missing historical endpoint proof.

\paragraph{Exact source estimates.}
Use the definitions of SN2, except that $n_s=\beta/q^2$ need not be
bounded. The exact sphere SN3 remains valid. Put $b=\beta/q$ and
$z=\kappa\beta/q^2$. Then
\[
\begin{gathered}
 h=q\tau\in\operatorname{Inf},\quad\theta\sim1,\quad
 e=\kappa\tau/\ell\in\operatorname{Inf},\quad M\sim1,\\
 b=\frac{\tau(B/S)-\ell H_0}{2\theta(P/S)},\qquad
 \mu\ell^2=\frac{\tau(1+x)(1+\sigma r)}{2M(P/S)},\\
 Jd_s+\theta z=\frac{\ell^2\mathcal H}{2(P/S)}\sim\tfrac12.
\end{gathered}\tag{BN2}
\]
In particular $b,z,d_s,\mu,\tau$ are bounded, $\tau$ and $\mu$
have positive real lower bounds, and at least one of $d_s,z$ is a
positive bounded unit. The exact inequality $\beta\ge0$ supplies
the lower bound on $\tau$. Also
$c_0,a_0,a_1\sim1/2$, $D,\alpha,m,\beta$ are infinitesimal.
Although the two terms below need not have small difference before
subtraction, their displayed bounded formulas in BN2 give
\[
                  A_{00}:=\mu\ell^2-b\sim\tfrac12.
 \tag{BN3}
\]
This argument takes residues only of bounded source coefficients and
never uses an estimate divided by $q$. It allows arbitrary relative
rates of $S,\kappa,q$.

\paragraph{A retained circuit and two actual positive rows.}
Use the full reciprocal parent SN5, the two unrestricted normals SN6,
and the shifted exact scalar dictionary SN9. Let $\zeta=\ell/\kappa$,
$a=u^2,c=w^2$, $x=\kappa S$, and assume $u,w\ge0$ by the original
sign gauges. Both may be zero. For every state define
\[
\begin{array}{ll}
 A_1=\sqrt\beta\,uZ,&A_2=\sqrt\beta\,W/(\zeta\sqrt q),\\
 A_3=\sqrt\beta\sqrt{x/q}\,wZ,
 &A_4=\sqrt\beta\,uwW/(\zeta q\sqrt x),
\end{array}\tag{BN4}
\]
\[
 \Psi=\sum_{i=1}^4 A_i^2-\frac{4\beta}{\zeta q}uwWZ,
 \qquad E_1=A_1-A_2,\quad E_2=A_3-A_4.
 \tag{BN5}
\]
Every monomial and its radical is a product of a pure-source and a
pure-testing element. This remains true if $x/q$ is unbounded.
The exact inequality meant here is a separated SOS certificate:
\[
                    \Psi-\tfrac12(E_1^2+E_2^2)\in\Sigma T^2.
 \tag{BN6}
\]
Indeed, if $WZ\ge0$, then $A_1A_2,A_3A_4\ge0$ and
\[
 \Psi=E_1^2+E_2^2+
              2(\sqrt{A_1A_2}-\sqrt{A_3A_4})^2.
\]
The radicals displayed here split into pure-source and pure-testing
radicals. If $WZ<0$, use instead
\[
 \Psi-\tfrac12(E_1^2+E_2^2)
 =\tfrac12(A_1+A_2)^2+\tfrac12(A_3+A_4)^2
                  +\frac{4\beta}{\zeta q}|uwWZ|.
\]
The last positive separated monomial is a square. Thus BN6 is valid
for all states, including $WZ=0$. At the actual source contact every
$A_i$ agrees, so $\Psi$ and its first contact derivatives vanish.

\paragraph{The exact scalar after circuit subtraction.}
Write $N_0=W/\zeta-uwZ=N/\ell$, where $N$ is SN7, and keep
$G,B,B_u,R$ from SN7. Thus
\[
 R=\frac{\sqrt S}{\zeta}G-uN_0.
 \tag{BN7}
\]
Subtract the entire circuit BN5 from the exact scalar graph in SN9.
Define
\[
\begin{gathered}
 A=\zeta^2 A_N-b,\qquad A_P'=A_P-zc/\ell^2,\\
 K'=K-2\beta/(\zeta q),\quad J'=K'-\ell A_P',
 \quad L'=J'-A/\zeta,\\
 H'=F_Z-\beta a-\beta xc/q-\ell^2ScA_P'
                          +ac(A-2\zeta J').
\end{gathered}\tag{BN8}
\]
Then $F_0-\Psi=A N_0^2+A_P'G^2-2\zeta uwZL'N_0+H'Z^2$.
All coefficients of $A$ and $A_P'$ are bounded:
\[
 A=\mu\ell^2(1+ra+ec)^2
       +(\beta+\tfrac12)\tau^2c(1+c)-b,
 \qquad \overline A=\overline{A_P'}=\tfrac12.
 \tag{BN9}
\]
Here and below the residue is taken at infinitesimal $u,w$.

The circuit subtraction has especially simple exact effects on the
SN10 quotient dictionary:
\[
\begin{gathered}
 L'_a=L_a,\qquad L'_c=L_c+\kappa n_s/\ell,\qquad
 H'_{20}=H_{20},\quad H'_{02}=H_{02},\quad
 H'_{11}=H_{11}-2n_s,\\
 H'=S\delta^2(d_s+cH_{02})
               +c(a-x)\delta H'_{11}+c(a-x)^2H_{20}.
\end{gathered}\tag{BN10}
\]
In particular no coefficient $n_s$ remains on its own in the final
conductors. With $C_R$ and $L_a$ from SN11--SN13 put
\[
 C_G'=\kappa(H_{11}-2n_s)/\ell,\qquad
 C_N'=-H_{11}-2\ell L_c/\kappa.
 \tag{BN11}
\]
These are bounded. The cancellation for the second coefficient is
explicit:
\[
 C_N'=-Sa_0\ell^2r+(\beta+m+\tfrac12)\tau^2(c+1/\ell)
                 +2\ell\kappa\mu\tau-m\tau^2.
 \tag{BN12}
\]
For $C_G'$ use the bounded formula SN13 and $\kappa n_s=z$.
The exact full scalar is consequently
\[
\begin{split}
 F_0={}&\Psi+A N_0^2+A_P'G^2+d_sB^2+H_{02}(wB)^2
       +C_RR^2-2\sqrt S\,L_aGR\\
 &+C_G'G(wB)+C_N'N_0(wB_u).
\end{split}\tag{BN13}
\]
This is an identity in unrestricted state coordinates, not a
restriction to a contact graph in the full parent.

\paragraph{Free helper rows across the whole source sphere.}
First reserve the positive rows $E_1,E_2$ from BN6. If $d_s$ is a
unit, $wB$ has infinitesimal cost on $\sqrt{d_s}B$. Otherwise
$z$ is a positive bounded unit. The exact identity
\[
 E_2=-\sqrt z\,\{(w/\ell)G+
                          \sqrt S\,w(w-\sqrt q)Z\}
\]
then gives
\[
 wB=-\frac{w+\sqrt q}{\sqrt z}E_2
                   -\frac{w(w+\sqrt q)}\ell G.
 \tag{BN14}
\]
Every coefficient in BN14 is infinitesimal, so any fixed real pivot
on $wB$ is free at infinitesimal Gram cost. Only the chosen pure-source
unit $\sqrt z$ is inverted. Literal $\beta=0$ uses the first branch.

After this pivot has been purchased, the normalized SN8 identity is
\[
 \delta^2N_0^2+\frac{2\kappa\delta}{\ell}G(wB)
          +2\kappa(wB)^2-(wB_u)^2
                         =(\delta W/\zeta)^2.
 \tag{BN15}
\]
It buys any fixed pivot on $wB_u$ at infinitesimal cost. The order of
these two purchases is deliberate; no weak $B$ diagonal is needed
in the second step.

\paragraph{Both full normals and their only additional row.}
Use the complete $H_X,H_Y,H_{XY},Y_P,Y_a,Y_c$ from SN18 and its
bounded $\widehat A_X,\widehat X_a,X_c$ from SN17. These formulas
only require $\tau,\mu,\beta$ bounded, not $n_s$ bounded.
The entire first-normal part is
\[
\begin{split}
 L_X={}&\ell w\widehat A_XN_0
       +\frac{uw\sqrt S}{\ell}\widehat X_aG
       +d_s\sqrt S\,quB-bB_u+wX_c(wB_u),\\
 L_Y={}&\sqrt{r(1-r)}\{Y_PG+\sqrt S\,Y_aR+Y_c(wB)\}.
\end{split}\tag{BN16}
\]
The apparent extra $B_u$ coupling is paid by the actual circuit row
$E_1$, using the exact, zero-face-safe identity
\[
 bB_u=-(w+\sqrt q)\sqrt b\,E_1-b(w+\sqrt q)N_0.
 \tag{BN17}
\]
All coefficients on the right have zero residue because $b$ is
bounded. No inverse of $b$ or $\beta$ is used. In the $B$ term in
$L_X$, write $d_sB=\sqrt{d_s}(\sqrt{d_s}B)$; this is likewise
face-safe when $D=0$. All direct normal coefficients are
infinitesimal, and $H_X,H_Y\sim1/2$, $H_{XY}$ is infinitesimal.
There is no missing $X$ or $Y$ normal and no imposed state ratio.

\paragraph{Allocation and all-tests conclusion.}
Keep a fixed positive fraction of each circuit row in BN6.
Use the strict base rows $(N_0,G,X_e,N_y,\sqrt{d_s}B,E_1,E_2)$,
omitting any zero row. Their main four diagonal residues are $1/2$;
the last three have fixed positive pivots wherever their rows are
nonzero. Terms containing $R$ are infinitesimal matrix perturbations
by BN7. Purchase sufficiently large fixed pivots on $wB,wB_u$ by
BN14--BN15, and use them for the bounded cross terms in BN13 and
BN16. Formula BN17 pays the remaining first-normal row at
infinitesimal cost. Every purchase is taken from an explicitly
retained positive square or from the exact SOS remainder BN6.
The resulting bounded Gram has strictly positive definite real
residue. The surviving strict bounded-Gram lemma therefore gives
finite tensor SOS for SN5 with both unrestricted normals restored.

This proves all infinitesimal $u,w$, including the literal axes.
For noninfinitesimal Alice use the unchanged adjacent-$C,A$ theorem
CE1--CE10; $C,A$ are unbounded on BN1. For the other Bob charts use
the surviving small-$e$ unmatched theorem SE10--SE12, since
$\overline\ell=1$. The sign gauges and real/imaginary state split
are the original exact ones. The local argument already handles
either sign of $WZ$, so no state normalization or state-chart
exclusion is required. Both literal $D=0$ and $\beta=0$ are covered
whenever they intersect BN1.

\paragraph{Verification status.}
The accompanying exact checker is
\texttt{check\_c144\_bounded\_h\_ratio.py}. Bounds and separated
SOS legality are the analytical argument above. At creation this
theorem is a candidate pending independent lead review. The source
case $h/q$ unbounded and the $\ell$-infinitesimal endpoint remain
separate; no whole-C14.4 conclusion is asserted here.
